\section{问题与离散模型}
一块单位方板有两处固定热源。预测给出每个格子的温度后，核对热量平衡比重算整张温度图便宜。若同一块板上到处都少算了流出的热量，这种误差会逐格积累；若相邻格子一正一负，热传导可能抵消部分误差。将残差压缩成一个最大绝对值，会抹去这种差别。本文研究的问题是：保留每个区域的残差极值及符号，能否获得严格受控、计算可行的全场误差上界？

\paragraph{研究对象。}
取单位正方形、导热系数 $k=1$、环境温升 $0$。两块热源分别为
$[0.2,0.4]^2$ 上的 $1$ 和 $[0.6,0.8]^2$ 上的 $2$，其余区域为 $0$。
四边 Robin 换热系数 $\mu=(h_L,h_R,h_D,h_U)\in[1/5,5]^4$。
用 $n\times n$ 个单元的中心有限体积法离散，$\Delta=1/n$，
$N=n^2$。每个内部面在相邻单元之间产生单位导热权重；
每个外边界面在相应对角元上增加 $\beta(h)=\Delta/(1/h+\Delta/2)$。
载荷 $b$ 由热源和单元的交叠面积积分得到，总热量为 $0.12$。
于是对于内部网格图的 Laplacian $K$ 和各边的对角标记矩阵 $D_s$，
\[
 A_\mu=K+\sum_s\beta(h_s)D_s,\qquad A_\mu T=b.
\]
角点单元累加两个外表面；单格网格累加四面。下文的真解均指这个有限维系统的精确解。若取近似场 $\widetilde T$，定义残差 $r=b-A_\mu\widetilde T$，误差为 $e=T-\widetilde T$，于是 $e=G_\mu r$，其中 $G_\mu=A_\mu^{-1}$。

\paragraph{与已有工作的关系。}
变 Robin 系数导热的降阶近似与误差控制已有系统研究\cite{boyaval2009}；
对单调矩阵或 $H$ 矩阵线性系统构造分量误差界也有成熟方法\cite{minamihata2013,minamihata2015}。
我们聚焦一个有限问题：对指定参数区间和给定块摘要，哪些误差保证能廉价地预先准备，哪些空间与符号信息确实减少了多余的报警？这是一项有适用范围的推导与可复核的算法研究。前期导热项目的求解器和冻结参数为本次对比提供共同起点。
